\subsection{Profiled Gaussian marginal likelihood in penalized spectral coordinates}
\label{sec:gaussian-likelihood}

\paragraph{Why introduce a probability model?}
The quadratic PLS objective \eqref{eq:pls} defines a fitted
trend for any prescribed $\lambda$ without assuming a
probability distribution. A \emph{likelihood}, by contrast,
requires explicit stochastic assumptions about measurement
errors and departures from the reference difference $\mu$.
The following model gives one such interpretation; it is
not a distribution-free consequence of minimizing PLS.
The construction belongs to the established literature
on restricted maximum likelihood (REML) and
generalized maximum likelihood (GML) for smoothers
\citep{patterson1971,wahba1985,wood2011}.
A directly related Gaussian marginal likelihood for
Whittaker--Henderson smoothing is derived by
\citet[Sec.~5.3]{biessy2026}.

Let $D_d=U\Sigma V^\top$ be the reduced SVD from
\eqref{eq:svd-target}, with $r=n-d$,
$\delta_j=\sigma_j^2>0$, and let
$W\in\R^{n\times d}$ have orthonormal columns
spanning $\ker(D_d)$. Thus $W^\top V=0$,
$V^\top V=I_r$, and the columns of $[W,V]$
form an orthonormal basis of $\R^n$. Let
$p_\mu$ be \emph{any} deterministic vector
satisfying $D_dp_\mu=\mu\mathbf1_r$.
Every possible latent trend can then be expressed
as
\begin{equation}\label{eq:likelihood-decomposition}
\tau=p_\mu+W\beta+V\theta,\qquad
\beta\in\R^d,\ \theta\in\R^r .
\end{equation}
Here $\beta$ represents the unpenalized polynomial
degrees of freedom; $D_dW=0$, so the penalty does
not determine $\beta$. The vector $\theta$ comprises
the $r$ penalized coordinates. We treat $\beta$
as an unknown \emph{fixed effect}, not as a
zero-mean Gaussian random variable.

\begin{theorem}[Gaussian contrast likelihood, REML and profiling]
\label{thm:reml}
Fix an integer $1\le d<n$ and a specified numeric
$\mu$. Suppose that conditional on the unknown fixed
vector $\beta$ the model is
\begin{equation}\label{eq:gaussian}
\begin{aligned}
y&=p_\mu+W\beta+V\theta+\varepsilon,\\
\varepsilon&\sim N_n(0,\sigma_\varepsilon^2I_n),\\
\theta&\sim N_r(0,\sigma_u^2\Sigma^{-2}),\\
\theta&\perp\!\!\!\perp\varepsilon,\qquad
\lambda=\frac{\sigma_\varepsilon^2}{\sigma_u^2}>0.
\end{aligned}
\end{equation}
Both variances are strictly positive.
The assumption on $\theta$ is equivalent to
$D_d\tau-\mu\mathbf1_r\sim
N_r(0,\sigma_u^2I_r)$.
The \emph{error contrasts}
\begin{equation}\label{eq:contrast-z}
z=V^\top(y-p_\mu)
\end{equation}
eliminate $\beta$ and have independent coordinates
\begin{equation}\label{eq:contrast-variance}
z_j\sim N\left(0,\sigma_\varepsilon^2
c_j(\lambda)\right),\qquad
c_j(\lambda)=1+\frac{1}{\lambda\delta_j}.
\end{equation}
For any observed $z\ne0$, profiling out
$\sigma_\varepsilon^2$ yields
\begin{equation}\label{eq:ml}
\mathcal L(\lambda)=
r\log\left(\frac1r\sum_{j=1}^r
\frac{z_j^2}{c_j(\lambda)}\right)
+\sum_{j=1}^r\log c_j(\lambda).
\end{equation}
Minimizing $\mathcal L$ over $\lambda>0$
is equivalent to maximizing the Gaussian
restricted (contrast) likelihood for this model.
The displayed criterion is profiled negative
twice log likelihood, up to a constant independent
of $\lambda$. Its statistical construction is
established prior art in restricted and
smoothing-parameter likelihood
\citep{patterson1971,wahba1985,wood2011,biessy2026}.
\end{theorem}

The contrast decomposition and profiling argument are proved
in Appendix~\ref{app:likelihood-proofs}.


\paragraph{Selecting the historical variance ratio.}
The PLS weight has the variance-component meaning
$\lambda=\sigma_\varepsilon^2/\sigma_u^2$
already present in the original stochastic
formulation \citep{guerrero2007,cortes2017}.
The operational contribution here is to select
that ratio using the \emph{established}
Gaussian restricted/marginal likelihood
in the same bounded coordinate used for the
trend fit. For a finite, interior optimum,
\begin{equation}\label{eq:variance-ratio-selection}
\widehat\lambda\in\arg\min_{\lambda>0}
\mathcal L(\lambda),\qquad
\widehat S=S_d(\widehat\lambda).
\end{equation}
Since $S_d$ is a bijection on the interior,
one can equivalently minimize
$\mathcal L(\Lambda_{n,d}(s))$ over
$s\in(0,1)$. For a finite evaluated set
$\mathcal S\subset(0,1)$, however, the
software returns only the smallest score
\emph{on that grid}; a global optimum is
not thereby certified. Once a finite positive
$\widehat\lambda$ is selected, the
conditional component estimates are
\begin{equation}\label{eq:variance-component-profiles}
\begin{aligned}
\widehat{\sigma}_\varepsilon^2
&=\frac{J_{\widehat\lambda}}{n-d},\\
\widehat{\sigma}_u^2
&=\widehat{\sigma}_\varepsilon^2/
\widehat\lambda.
\end{aligned}
\end{equation}
These estimates depend on the Gaussian
assumptions and on a \emph{prescribed} $\mu$.
At a boundary optimum, a vanishing variance
component or other limit must be handled
separately; \eqref{eq:variance-component-profiles}
is an interior result. We therefore
contribute an explicit, reproducible
variance-ratio selection route within the
normalized PLS framework, \emph{not} a
new likelihood principle.

\paragraph{What the likelihood does and does not estimate.}
The quantity $z_j$ is not the observation $y_j$:
it is a projection of all $n$ observations onto
a right singular vector of $D_d$. The term
$\sigma_\varepsilon^2$ measures the assumed
observation-noise variance; $\sigma_u^2$ measures
random departures of order-$d$ differences from
the \emph{fixed} reference $\mu$; and
$\lambda=\sigma_\varepsilon^2/\sigma_u^2$
is their variance ratio. The scalar $c_j$
is a \emph{variance multiplier}, not a
variance ratio by itself. Larger $\lambda$
means stronger smoothing because the difference
innovation variance is smaller relative to
the observation noise variance. The $d$
polynomial directions disappear from $z$
because their coefficients are treated as
unpenalized nuisance fixed effects, which
is why the profile variance divisor is
$r=n-d$, not $n$.

Neither maximizing this likelihood nor
minimizing PLS automatically gives the same
$\lambda$ as CV or GCV; they solve distinct
statistical selection problems. The case $z=0$
has a zero-variance boundary degeneracy
under profiling and is excluded from the
interior formula in Theorem~\ref{thm:reml}.
At $\lambda=0$ the covariance parametrization
in \eqref{eq:contrast-variance} is only a
limit, since it would require unbounded
$\sigma_u^2$ for fixed positive
$\sigma_\varepsilon^2$.

\begin{proposition}[Penalized quadratic and determinant identities]
\label{prop:det}
For $\lambda>0$ let
$B_d=D_dD_d^\top$ and let
$\widehat\tau$ solve \eqref{eq:rhs}. Define
\begin{equation}\label{eq:likelihood-J}
J_\lambda=
\|y-\widehat\tau\|_2^2+
\lambda\|D_d\widehat\tau-\mu\mathbf1_r\|_2^2.
\end{equation}
Then
\begin{equation}\label{eq:quadraticidentity}
J_\lambda=\sum_{j=1}^r
\frac{z_j^2}{c_j(\lambda)}
\end{equation}
and
\begin{equation}\label{eq:detidentity}
\sum_{j=1}^r\log c_j(\lambda)
=\log\det A_\lambda-r\log\lambda
-\log\det B_d.
\end{equation}
Thus the \emph{same} profiled likelihood
can be computed using structured Cholesky
factorizations, without forming a dense
inverse or using an SVD at every candidate.
\end{proposition}
The orthogonal quadratic minimization and log-determinant
argument are given in Appendix~\ref{app:likelihood-proofs}.


\paragraph{Precisely what is prior art?}
The contrast/restricted-likelihood principle
is classical \citep{patterson1971}; generalized
maximum likelihood for smoothing-parameter
selection predates this article
\citep{wahba1985}, as do stable REML and
marginal-likelihood algorithms for penalized
smoothers \citep{wood2011}. In the
Gaussian original Whittaker--Henderson
model, \citet[Sec.~5.3]{biessy2026}
gives a marginal likelihood with a minimized
penalized quadratic and a
log-determinant correction. On specialization
to the unweighted model with $\mu=0$,
its quadratic and determinant are
\eqref{eq:quadraticidentity} and
\eqref{eq:detidentity}, up to
the overall likelihood sign convention
and constants independent of $\lambda$.
Our expression in \eqref{eq:ml}
is an \emph{algebraic representation of
existing Gaussian marginal/REML methodology
in the coordinates of $D_d$}, not a
newly invented likelihood. The mathematical
contribution claimed here is the explicit
integration with fixed-$\mu$ reconstruction,
the normalized smoothness coordinate,
and complete algorithmic workload accounting;
its novelty relative to existing computational
treatments still requires independent review.

\paragraph{Fixed versus estimated reference mean.}
Everything above conditions on a
\emph{prespecified numeric} $\mu$.
Substituting
$\widehat\mu=\mathbf1_r^\top D_dy/r$
changes the observation-to-fit map
and the statistical estimation problem.
The displayed $r$-dimensional contrast
likelihood cannot then be called
joint unknown-$\mu$ REML without
deriving the altered fixed-effect
space, degrees of freedom, and covariance.
If $\mu$ is freely estimated as a
new regression coefficient, an additional
nuisance direction must be accounted for.
Likewise, fixed-$\mu$ conditional
LOOCV is not automatic foldwise refitting
of an estimated reference mean.

\paragraph{Computing contrasts without a high-degree shift.}
Although $z=V^\top(y-p_\mu)$ is an
unambiguous \emph{mathematical} definition,
forming $p_\mu$ explicitly may be numerically
unwise for high $d$ or large $n$.
The SVD yields the equivalent identity
\begin{equation}\label{eq:contrast-shiftfree}
z=V^\top y-\mu\Sigma^{-1}
U^\top\mathbf1_r.
\end{equation}
This avoids explicitly constructing
and subtracting the degree-$d$
polynomial $p_\mu$, but division by
a very small $\sigma_j$ can still
amplify numerical errors. Finite-precision
residual checks remain necessary
(Section~\ref{sec:stability}).
